\documentclass[11pt]{article}
\usepackage[margin=1in]{geometry}
\usepackage{enumitem}
\setlist[itemize]{leftmargin=*, itemsep=2pt, topsep=2pt}

\title{LearnByInterrogation}
\author{}
\date{}

\begin{document}
\maketitle

\section*{Inspiration}
Notes apps store knowledge; they don't test it. We wanted the Socratic method — the oldest, most rigorous teaching tool there is — applied automatically to whatever a learner writes, so understanding is verified, not assumed.

\section*{What it does}
\begin{itemize}
  \item User writes short notes (formulas, definitions, claims) tagged by type.
  \item A local, heuristic engine turns notes into a knowledge graph of concepts and links.
  \item The graph surfaces gaps: undefined terms, isolated ideas, unconnected-but-related concepts.
  \item An AI ``Socrates'' interrogates the user about each new note and its gaps — never explains, only questions.
  \item Split view: graph on the left, Socratic dialogue on the right.
\end{itemize}

\section*{How we built it}
\begin{itemize}
  \item React + TypeScript + Tailwind, deployed via GitHub Actions to GitHub Pages.
  \item \texttt{react-flow} for an interactive, draggable knowledge graph.
  \item Client-side embeddings (\texttt{transformers.js}, MiniLM) in a Web Worker — no server, no upload.
  \item Concept/gap extraction via local heuristics (keyword/pattern matching + embedding similarity), not an LLM call.
  \item Interrogation powered by Groq (BYOK), retrieving relevant notes as context (RAG) for each question.
  \item All state — notes, embeddings, graph, chat — persisted in IndexedDB.
\end{itemize}

\section*{Challenges we ran into}
\begin{itemize}
  \item Running a real embedding model fully in-browser without a backend.
  \item Building a useful knowledge graph without calling an LLM for every note.
  \item Making the graph feel alive (auto-layout, drag-to-rearrange) without fighting React's re-render cycle.
  \item Prompting an LLM to \emph{ask} rather than \emph{answer}.
\end{itemize}

\section*{Accomplishments that we're proud of}
\begin{itemize}
  \item A fully local-first architecture: notes never leave the browser except the interrogation call itself.
  \item Automatic gap detection that reliably surfaces genuinely undefined or disconnected ideas.
  \item Interrogation that starts itself the moment you write something — zero friction.
\end{itemize}

\section*{What we learned}
\begin{itemize}
  \item Small, local models are viable for real product features, not just demos.
  \item Constraining an LLM to ask questions (not answer) requires very deliberate prompting.
  \item Simple heuristics can substitute for expensive AI calls in graph construction.
\end{itemize}

\section*{What's next for LearnByInterrogation}
\begin{itemize}
  \item Spaced-repetition scheduling driven by graph gaps.
  \item Multi-session memory of what the user has and hasn't yet understood.
  \item Import from existing notes (PDF, Markdown, Notion).
  \item Shared/collaborative graphs for study groups.
\end{itemize}

\end{document}
